\section{激活函数}

\subsection{为什么需要激活函数}

激活函数的唯一目的是引入\textbf{非线性}。没有激活函数，无论多少层线性运算（卷积或全连接）的堆叠，整个网络仍然等价于一个线性变换。

激活函数的放置位置：通常在\textbf{卷积层之后}和\textbf{全连接层之后}。

\subsection{Sigmoid：历史上的经典}

\begin{figure}[H]
\centering
\includegraphics[width=0.85\textwidth]{slides-images/slide-014.jpg}
\caption{Sigmoid 函数：$\sigma(x) = \frac{1}{1 + e^{-x}}$，输出范围 $(0, 1)$\protect\footnotemark}
\end{figure}
\footnotetext{Slides 第15页。}

$$\sigma(x) = \frac{1}{1 + e^{-x}}$$

Sigmoid 将任意实数映射到 $(0, 1)$ 区间。它在深度学习早期非常流行，但有一个\textbf{致命缺陷}：

\begin{warningbox}{Sigmoid 的梯度消失问题}
Sigmoid 的导数 $\sigma'(x) = \sigma(x)(1 - \sigma(x))$ 的最大值仅为 0.25（在 $x = 0$ 处）。在 $|x|$ 较大的区域，梯度接近于零。

这意味着：当输入值很大或很小时（这在深层网络中经常发生），梯度几乎为零，参数无法更新。经过多层 sigmoid 后，靠近输入层的梯度会指数级衰减——这就是\textbf{梯度消失}（vanishing gradient）问题，曾经严重制约了深度网络的训练。
\end{warningbox}

\subsection{ReLU：简单而有效}

\begin{figure}[H]
\centering
\includegraphics[width=0.85\textwidth]{slides-images/slide-015.jpg}
\caption{ReLU 函数：$f(x) = \max(0, x)$\protect\footnotemark}
\end{figure}
\footnotetext{Slides 第16页。}

$$f(x) = \max(0, x)$$

ReLU（Rectified Linear Unit）的出现是深度学习发展的关键里程碑：
\begin{itemize}
    \item 正半轴梯度恒为 1——\textbf{不会梯度消失}
    \item 计算极其简单——只需一次比较操作
    \item 引入了恰好足够的非线性
\end{itemize}

ReLU 的缺点：负半轴梯度为零（``死神经元''问题），一旦某个神经元的输入长期为负，它将永远不再激活。

\subsection{GELU 与 SiLU：现代选择}

\begin{figure}[H]
\centering
\includegraphics[width=0.85\textwidth]{slides-images/slide-017.jpg}
\caption{GELU 函数及其与 ReLU 的对比\protect\footnotemark}
\end{figure}
\footnotetext{Slides 第18页。}

\textbf{GELU}（Gaussian Error Linear Unit）：

$$\text{GELU}(x) = x \cdot \Phi(x)$$

其中 $\Phi(x)$ 是标准正态分布的累积分布函数。

\textbf{SiLU / Swish}：

$$\text{SiLU}(x) = x \cdot \sigma(x)$$

其中 $\sigma$ 是 sigmoid 函数。

\begin{importantbox}{现代激活函数的设计思路}
GELU 和 SiLU 的共同特点：
\begin{itemize}
    \item 在 $x \to +\infty$ 时趋近于 $f(x) = x$（线性，梯度为 1）
    \item 在 $x \to -\infty$ 时趋近于 $f(x) = 0$（类似 ReLU）
    \item 在 $x \approx 0$ 附近是\textbf{平滑的}（不像 ReLU 有尖角）
    \item 在负半轴有\textbf{非零梯度}（避免死神经元问题）
\end{itemize}
GELU 是当前 Transformer 中使用最广泛的激活函数。
\end{importantbox}

\begin{figure}[H]
\centering
\includegraphics[width=0.85\textwidth]{slides-images/slide-018.jpg}
\caption{常见激活函数对比：Sigmoid、ReLU、GELU、SiLU 等\protect\footnotemark}
\end{figure}
\footnotetext{Slides 第19页。}

\begin{knowledgebox}{激活函数的选择建议}
\begin{itemize}
    \item \textbf{传统 CNN}：ReLU（简单高效）
    \item \textbf{Transformer}：GELU（平滑且有非零负半轴梯度）
    \item \textbf{现代 CNN}：SiLU / Swish（与 GELU 性能接近）
    \item \textbf{不要使用}：Sigmoid 或 Tanh 作为隐藏层激活（梯度消失风险）
\end{itemize}
\end{knowledgebox}

\subsection{激活函数的放置位置}

\begin{figure}[H]
\centering
\includegraphics[width=0.85\textwidth]{slides-images/slide-019.jpg}
\caption{激活函数的放置：在线性层（全连接或卷积）之后\protect\footnotemark}
\end{figure}
\footnotetext{Slides 第20页。}

激活函数总是放在线性运算之后：
\begin{itemize}
    \item 全连接层 $\rightarrow$ 激活函数
    \item 卷积层 $\rightarrow$ 激活函数
    \item 归一化层可以放在激活函数之前或之后（具体取决于架构设计）
\end{itemize}

\subsection{本章小结}

激活函数是深度网络表达能力的关键。从 Sigmoid（梯度消失严重）到 ReLU（简单高效但有死神经元），再到 GELU（平滑且处处可微），激活函数的发展推动了深度学习的进步。选择合适的激活函数对训练稳定性和最终性能都有显著影响。

%% ========== Part 5: CNN 经典架构 ==========

